\documentclass[11pt]{amsart}
\usepackage{geometry}                % See geometry.pdf to learn the layout options. There are lots.
\geometry{letterpaper}                   % ... or a4paper or a5paper or ... 
%\geometry{landscape}                % Activate for for rotated page geometry
%\usepackage[parfill]{parskip}    % Activate to begin paragraphs with an empty line rather than an indent
\usepackage{graphicx}
\usepackage{amssymb}
\usepackage{epstopdf}
\DeclareGraphicsRule{.tif}{png}{.png}{`convert #1 `dirname #1`/`basename #1 .tif`.png}

\title{PS 4.35}
%\date{}                                           % Activate to display a given date or no date

\begin{document}
\maketitle
%\section{}
%\subsection{}
\noindent Simule por medio de un diagrama de flujo un directorio telefónico que lea nom­ bres y números telefónicos en dos arreglos paralelos. Los nombres están ordena­ dos alfabéticamente en orden ascendente.
\newline

Cuando el usuario teclee un nombre, el programa debe verificar si éste ya es­ taba en el directorio, en cuyo caso imprimirá el número telefónico. De no ser así y si hay capacidad, pedirá al usuario que escriba el teléfono y agregará nombre y número a los arreglos.
\newline

\noindent Datos: NOMBRE [1..N], TELEFONO [1..N] (arreglos unidimensionales de tipo cadena de caracteres, 1 $\leq$ N $\leq$ 1000).


\end{document}  